\scanpage{54}
of a simplicial object (space) --- \uline{over}
$\Gamma(\varnothing)=\Sigma^{<}(0)$. Here it is a matter of a relative
\margin{NB $\Gamma(n)=\Sigma^{*}(n+1)$}
simplicial space, and the conditions $b_1,b_2,b_3$
concern this structure --- they say
essentially that this structure is associated
with a local cylindrical stratification of
$\Sigma^{<}(0)=\Sigma$, and ``ordered'' (i.e. the set of
strata of codimension $1$ passing through an $x\in\Sigma$ are
\deleted{\unsure{is}} \unsure{ordered|ordered}, and this order relation varies
continuously with $x$ in an obvious sense).

The \deleted{$\Sigma^{<}(d)$ \unsure{are identified|are identified}} $\Gamma(d)$\inserted{$=\Sigma^{<}(d+1)$} are then identified
with the \deleted{parts} spaces of the \inserted{(germs of)} strata of codimension
$d+1$, \deleted{the} \inserted{germs} identified \deleted{(thanks to the order
structure)} with the sets of germs of strata
of codimension $1$ \deleted{\unsure{of}} \inserted{\unsure{at an}} $=$ $x$. The semi-simplicial map\unsure{s}
between the $\Gamma(d)$ then arise
from the order structure --- \deleted{\unsure{there one}}
\unsure{to|can} reconstruct the flag spaces
associated with the stratification, using the
\deleted{\unsure{this|these}} $\Gamma(d)$, thanks to this ordered structure.
In summary, the datum of $\Gamma:d\mathrel{\mid\! -\!}\Gamma(d)$
\deleted{\unsure{i.e. of}} (i.e. of $\Sigma^{<}$ with the cofinal \unsure{semi|\deleted{semi}} simplicial \unsure{maps|maps})
is equivalent to that of a $\Sigma=\Sigma^{<}(0)$,
\margin{53}
